% from ega4-16-fr.tex
\oldpage[IV-4]{8}
Soient $X$ un préschéma, $j:Y\to X$ une immersion localement de présentation
finie, $y$ un point de $Y$. Les conditions suivantes sont équivalentes:
\begin{enumerate}
\item[(a)] Il existe un voisinage ouvert $U$ de $y$ dans $Y$ tel que $j|U$ soit
un homéomorphisme de $U$ sur un ouvert de $X$.
\item[(b)] Il existe un entier $n>0$ tel que l'homomorphisme canonique
\[
  (\varphi_{n,n-1})_y:
  \mathcal{O}_{Y^{(n)},y}\longrightarrow\mathcal{O}_{Y^{(n-1)},y}
\]
soit bijectif.
\item[(c)] Il existe un entier $n>0$ tel que
$(\mathscr{G}r_n(j))_y=0$.
\end{enumerate}
De plus, si l'entier $n$ vérifie b) ou c), il existe un voisinage $V$ de $y$
dans $Y$ tel que $\mathscr{G}r_m(j)|V=0$ pour $m\geq n$ et que
$\varphi_{nm}|V:\mathcal{O}_{Y^{(m)}}|V\to\mathcal{O}_{Y^{(n)}}|V$ soit
bijectif pour $m\geq n$.
\end{proposition}

La question étant locale sur $Y$, on peut se borner au cas où $j$ est une
immersion fermée, $Y$ étant défini par un Idéal quasi-cohérent de type fini
$\mathscr{I}$ de $\mathcal{O}_X$. L'équivalence de b) et c), pour un $n$ donné,
est alors immédiate; en outre, comme
$\mathscr{I}^n/\mathscr{I}^{n+1}$ est un $\mathcal{O}_X$-Module de type fini,
il existe un voisinage ouvert $U$ de $y$ dans $X$ tel que
$\mathscr{I}^n|U=\mathscr{I}^{n+1}|U$ (0\textsubscript{I}, 5.2.2), donc aussi
$\mathscr{I}^n|U=\mathscr{I}^m|U$ pour $m\geq n$, ce qui prouve les dernières
assertions. Pour prouver que a) entraîne b), on peut se borner au cas où
l'espace sous-jacent à $Y$ est égal à l'espace sous-jacent à $X$, et où
$\mathscr{I}$ est engendré par un nombre fini de ses sections au-dessus de $X$:
comme $\mathscr{I}$ est alors contenu dans le Nilradical $\mathscr{N}$ de
$\mathcal{O}_X$ (I, 5.1.2), il est nilpotent, ce qui prouve b). Enfin, pour
prouver que b) entraîne a), on peut aussi se borner au cas où
$\mathscr{I}^n=\mathscr{I}^{n+1}$; alors, pour tout $y\in Y$, comme
$\mathscr{I}_y\subset\mathfrak{m}_y$, idéal maximal de
$\mathcal{O}_{X,y}$, on a nécessairement $\mathscr{I}_y^n=0$ en vertu du
lemme de Nakayama, puisque $\mathscr{I}_y$ est un idéal de type fini.
L'ensemble des $x\in X$ tels que $\mathscr{I}_x^n=0$ est donc un ouvert $U$ de
$X$ contenant $Y$ (0\textsubscript{I}, 5.2.2); comme d'autre part
$\mathscr{I}_x\neq 0$ pour $x\notin Y$, on a nécessairement $U=Y$.

\begin{corollary}[16.1.10]
\label{IV.16.1.10}
Pour que la restriction de l'immersion $j$ à un voisinage de $y$ dans $Y$ soit
une immersion ouverte (autrement dit pour que $j$ soit un isomorphisme local au
point $y$), il faut et il suffit que
$(\mathscr{G}r_1(j))_y=(\mathscr{N}_{Y/X})_y=0$.
\end{corollary}

La condition est évidemment nécessaire, et le raisonnement précédent, appliqué
pour $n=1$, prouve qu'elle est suffisante.

\begin{remarks}[16.1.11]
\EGAReaderTarget{ega4.c04.rem.16-1-11}{IV.16.1.11}
\begin{enumerate}
\item[(i)] Sous les conditions de la définition (16.1.1), la limite projective
du système projectif
$(\mathcal{O}_{Y^{(n)}},\varphi_{nm})$ de faisceaux d'anneaux sur $Y$ est
appelé l'invariant normal d'ordre infini de $f$, et parfois noté
$\mathcal{O}_{Y^{(\infty)}}$. Lorsque $X$ est un préschéma localement
noethérien, $j:Y\to X$ une immersion fermée, $Y$ étant donc un sous-préschéma
fermé de $X$ défini par un Idéal cohérent $\mathscr{I}$,
$\mathcal{O}_{Y^{(\infty)}}$ n'est autre que le complété formel de
$\mathcal{O}_X$ le long de $Y$ (I, 10.8.4), et
$Y^{(\infty)}=(Y,\mathcal{O}_{Y^{(\infty)}})$ le préschéma formel complété de
$X$ le long de $Y$ (I, 10.8.5). Dans tous les cas, on pourra dire que
$Y^{(\infty)}$ est le voisinage formel de $Y$ dans $X$ (pour le morphisme
$f$). Dans le cas particulier qu'on vient d'envisager, c'est donc le préschéma
formel limite inductive des voisinages infinitésimaux d'ordre $n$.
\item[(ii)] On notera que pour un morphisme de préschémas
$f=(\psi,\theta):Y\to X$, il peut se faire que l'homomorphisme
$\theta^\#:\psi^*(\mathcal{O}_X)\to\mathcal{O}_Y$ soit surjectif sans que $f$
soit une immersion
